% TOP_PROBLEM: {"id": 2303011, "title": "Research Problems in Function Theory — Problem 3.11", "category_id": 8, "category": {"id": 8, "name": "analysis", "display_name": "Analysis", "description": "Limits, continuity, calculus, and function theory.", "slug": "analysis", "order_index": 8, "created_at": "2026-07-31T15:26:25.671Z"}, "status": "open"}
% FIRST_POSTED: null
% ENTRY_KIND: statement_only
% Imported from: batch12.tex; UnsolvedMath ID 2303011
% Compile twice: lualatex 2303011.tex
\documentclass[11pt]{article}
\usepackage{fontspec}
\setmainfont{Latin Modern Roman}
\usepackage{amsmath,amssymb,mathtools,unicode-math}
\setmathfont{Latin Modern Math}
\usepackage{xurl,hyperref}
\hypersetup{colorlinks=true,urlcolor=blue,pdftitle={UnsolvedMath Problem 2303011}}
\newcommand{\sref}[2]{\hyperlink{source:#1:#2}{[S#2]}}
\newcommand{\cataloguescope}[1]{\paragraph{Mathematical question.}#1}
\begin{document}
% UnsolvedMath ID 2303011; source catalogue code AMR-022-3011
\setcounter{section}{2303010}
\section{Research Problems in Function Theory — Problem 3.11}\label{2303011}
{\small\sffamily UnsolvedMath ID 2303011 · Analysis}
\subsection{Definitions and mathematical statement}

\paragraph{Shared notation.}$\mathbb N=\{1,2,\ldots\}$, and $\mathbb Z,\mathbb Q,\mathbb R,\mathbb C$ denote the integers, rationals, reals and complex numbers. Any other conventions are specified locally.


A subharmonic function is an upper semicontinuous function with values in $[-\infty,\infty)$, not identically $-\infty$ on a connected component, that satisfies the mean-value inequality on every closed ball in its domain. A real harmonic function is a twice continuously differentiable function $u$ satisfying $\Delta u=0$. For $m\geq2$ and $n\geq1$, let $\mathcal P_{n,m}$ be the nonzero real homogeneous harmonic polynomials of degree $n$ on $\mathbb R^m$. Homogeneity means $u(tx)=t^n u(x)$. Put $A(r,u)=\inf_{|x|=r}u(x)$ and $B(r,u)=\sup_{|x|=r}u(x)$. The zero spherical mean gives $A(r,u)<0<B(r,u)$ when $u(0)=0$ and $u$ is nonconstant. For homogeneous polynomials, $-A(r,u)/B(r,u)$ is independent of $r>0$. A set $E\subset[1,\infty)$ has finite logarithmic measure when $\int_E dr/r<\infty$.
The source update attributes an extremal-ratio solution for homogeneous harmonic polynomials in dimensions $m\geq3$ to Armitage.
\paragraph{Statement recorded in the supplied source.}
Determine the optimal lower and upper bounds
\[L_{n,m}=\inf_{u\in\mathcal P_{n,m}}\frac{-A(1,u)}{B(1,u)},\qquad
U_{n,m}=\sup_{u\in\mathcal P_{n,m}}\frac{-A(1,u)}{B(1,u)}.\]
For odd $n$ both bounds are $1$. Also determine the sharp comparison of $-A(r,u)$ and $B(r,u)$ for nonpolynomial harmonic $u$ on $\mathbb R^m$ with $u(0)=0$: can the source's bound
\[-A(r,u)<B(r,u)\bigl(\log B(r,u)\bigr)^{m-1+\varepsilon}\]
be improved? Here $\varepsilon>0$ is arbitrary, the bound applies for all sufficiently large $r$ outside a set of finite logarithmic measure that may depend on $u$ and $\varepsilon$, and the displayed logarithm is taken where $B(r,u)>1$. Include the planar comparison $A(r,u)\sim-B(r,u)$ outside such a set as the known benchmark. \sref{2303011}{2}
\subsection{Short English statement}
Find the sharp extrema ratio for homogeneous harmonic polynomials and the sharp growth comparison between negative and positive extrema for nonpolynomial harmonic functions, allowing an exceptional set of finite logarithmic measure.
\subsection{Sources}
\begin{description}
\item[\hypertarget{source:2303011:1}{S1}] ULAM.AI and the UnsolvedMath Contributors, UnsolvedMath: A Curated Collection of Open Mathematics Problems, record 2303011 (AMR-022-3011). CC BY 4.0. \url{https://huggingface.co/datasets/ulamai/UnsolvedMath}.
\item[\hypertarget{source:2303011:2}{S2}] W. K. Hayman and E. F. Lingham, \emph{Research Problems in Function Theory}, arXiv:1809.07200v2 (2018), Problem 3.11 and Update 3.11. \url{https://arxiv.org/abs/1809.07200}.
\end{description}
\label{end:2303011}
\end{document}
